%%=============================================================================
%% ARTIGO -- Trabalho Final APPD (IDP, 2026)
%% Adaptacao do TemplateLatexIDP (commit 717f75f) para o formato de artigo,
%% conforme PLANO.md, secao 8. A classe idp.cls NAO e alterada; toda a
%% adaptacao (margens de 2,5 cm, sem pre-textuais de dissertacao) fica aqui.
%%
%% Compilar:  latexmk -xelatex main-artigo.tex   (ler main-artigo.log depois)
%%=============================================================================
% Tipo compartilhado pelo trabalho e pelo termo avulso.
% Valores: bacharelado | especializacao | mestrado | doutorado.
\def\NivelTrabalho{mestrado}

\documentclass[\NivelTrabalho,secaocontinua,links]{idp}

\usepackage{booktabs}
\usepackage{siunitx}
\usepackage{amsmath}
\graphicspath{{../outputs/figuras/}}

\addbibresource{referencias.bib}
%%=============================================================================
%% DADOS DO ARTIGO -- FONTE UNICA DA VERDADE (PLANO.md, secao 8.3)
%% Campos de banca, ficha catalografica e defesa NAO se aplicam ao artigo e
%% nao aparecem no PDF (main-artigo.tex nao chama \capa, \folhaderosto,
%% \fichacatalografica nem \folhadeaprovacao).
%%=============================================================================

%% --- Identificacao ---------------------------------------------------------
%% Titulo provisorio (politica, desfecho, lugar, periodo) -- revisar.
\titulo{Efeito do Programa Academia da Saúde sobre internações por doenças hipertensivas, AVC e diabetes nos municípios brasileiros, 2008--2024}
\subtitulo{}
\titulotraduzido{}

%% Nomes completos dos integrantes (separados por ponto e virgula).
\autor{}                                         % PREENCHER
\matricula{}

\orientador{Prof. Dr. Hélio Bomfim de Macêdo Filho}  % professor da disciplina
\coorientador{}

%% --- Vinculo institucional -------------------------------------------------
\curso{}
\programa{Mestrado Profissional em Administração Pública}
\area{}
\linhapesquisa{}

\cidade{Brasília}
\local{Brasília-DF}
\data{2026}

\datadefesa{}

%% --- Resumo ----------------------------------------------------------------
\palavraschave{Programa Academia da Saúde; avaliação de impacto; diferenças-em-diferenças; escore de propensão; internações hospitalares.}
\keywords{}
   % titulo, programa: fonte unica dos metadados

\begin{document}
\newgeometry{margin=2.5cm}   % margens exigidas para o artigo (orientacoes, secao 5)
\pagestyle{idp}\pagenumbering{arabic}

%% ---- Bloco de titulo do artigo (substitui capa e folha de rosto) ----------
\begin{center}
  {\large\bfseries \TituloTrabalho}\\[6pt]
  % PREENCHER: nomes completos dos integrantes, separados por ponto e virgula
  \NomeAutor\\[3pt]
  Avaliação de Políticas Públicas com Dados --- \NomePrograma, \NomeInstituicao\\
  \LocalTrabalho, \AnoTrabalho
\end{center}

%% RESUMO DO ARTIGO -- ate 250 palavras (orientacoes, secao 5), paragrafo unico.
%% Deve conter: pergunta, dados, metodos, resultado principal com magnitude e
%% IC 95%, conclusao. Escrever por ultimo (PLANO.md, Fase 7). Palavras-chave
%% ficam em config/dados.tex.
\begin{resumo}
[PREENCHER: resumo redigido pelos autores ao final da Fase 7.]
\end{resumo}
     % ate 250 palavras + palavras-chave
% %% RESUMO EM LINGUA ESTRANGEIRA (obrigatorio para mestrado/doutorado)
%% Traducao fiel do resumo em portugues. As keywords sao definidas em
%% config/dados.tex.

\begin{abstract}
This space is reserved to present the abstract of your work in a foreign
language; English is the suggested choice. The abstract should contain between
150 and 500 words, in Times New Roman size 12, single-spaced and justified. A
good abstract states the purpose of the overall work, the key concepts and
theoretical authors relied upon, the main steps and actions taken to implement
the proposal, the population involved, how the data were collected and the
period in which they were obtained. It should also present the main results and
the conclusions reached. The abstract must be written as a single paragraph,
with no first-line indentation, using the active voice and the third person
singular. Avoid symbols, formulas, equations and diagrams that are not strictly
necessary, as well as bibliographic citations.
\end{abstract}
 % opcional

\textual
%% Introdução -- conteúdo exigido (PLANO.md, seção 2.4):
%% Política e mecanismo; literatura (no mínimo 3 referências verificadas); pergunta causal; estimando (ATT); unidade de análise (município × ano).
\section{Introdução}
\label{sec:introducao}

% PREENCHER: redação dos autores.

%% Dados -- conteúdo exigido (PLANO.md, seção 2.4):
%% Fontes com link e data de acesso (CNES, SIH/SUS, IBGE); período; cobertura; construção de D, Y e X; decisões de limpeza com contagens; descritivas (tab01, fig02, fig03).
\section{Dados}
\label{sec:dados}

% PREENCHER: redação dos autores.

%% Estratégia empírica -- conteúdo exigido (PLANO.md, seção 2.4):
%% DAG (fig01); hipóteses de identificação (ignorabilidade e suporte comum; tendências paralelas e não antecipação); equações; erros-padrão clusterizados por município; diagnósticos planejados.
\section{Estratégia empírica}
\label{sec:estrategia-empirica}

% PREENCHER: redação dos autores.

%% Resultados -- conteúdo exigido (PLANO.md, seção 2.4):
%% Tabelas com estimativa, EP e IC 95\% (tab03, tab04, tab08); love plot (fig05); estudo de eventos (fig06, fig08). Números somente via \texttt{\textbackslash input} de outputs/tabelas.
\section{Resultados}
\label{sec:resultados}

% PREENCHER: redação dos autores.

%% Robustez -- conteúdo exigido (PLANO.md, seção 2.4):
%% Especificações R1--R6; placebos (tab05); robustez (tab06); poder e MDE (tab07).
\section{Robustez}
\label{sec:robustez}

% PREENCHER: redação dos autores.

%% Discussão e limitações -- conteúdo exigido (PLANO.md, seção 2.4):
%% O que se pode e o que não se pode afirmar; SUTVA (transbordamento); erro de medida no tratamento; ameaças à identificação. TEXTO DOS AUTORES.
\section{Discussão e limitações}
\label{sec:discussao-limitacoes}

% PREENCHER: redação dos autores.

%% Conclusão -- conteúdo exigido (PLANO.md, seção 2.4):
%% Resposta direta à pergunta, com implicações para a política. TEXTO DOS AUTORES (CLAUDE.md, regra 6).
\section{Conclusão}
\label{sec:conclusao}

% PREENCHER: redação dos autores.

%% Declaracao de uso de IA -- Anexo II da Politica de Uso de IA do IDP (2026).
%% Preencher na Fase 8 com modelo/versao exatos e datas reais (PLANO.md, 3.2e),
%% a partir de docs/ia/registro-uso-ia.md.
\section*{Declaração de uso de inteligência artificial}
\addcontentsline{toc}{section}{Declaração de uso de inteligência artificial}

Durante a preparação deste trabalho, os autores utilizaram o Claude Code
(Anthropic), modelo [PREENCHER: nome e versão exatos], entre
[PREENCHER: data inicial] e [PREENCHER: data final], para: (i) discutir o
desenho de pesquisa; (ii) apoiar a busca e o mapeamento de literatura, com
todas as referências verificadas nas fontes originais; (iii) gerar e depurar
código Python de download, extração, limpeza, construção da base, estimação e
visualização, nas etapas [PREENCHER: listar]; e (iv) revisar a formatação
LaTeX do artigo. Após o uso dessa ferramenta, os autores revisaram e editaram
o conteúdo em conformidade com o método científico, reexecutaram
integralmente os cadernos e assumem responsabilidade pelo conteúdo do
trabalho. Os prompts e as saídas relevantes estão disponíveis em
\texttt{docs/ia/} no repositório.
  % Anexo II da Politica de IA, antes das referencias

\postextual
\referencias
%% APENDICES (opcional): dicionario de variaveis, tabelas extras, link do repositorio.
\apendice{Repositório e reprodutibilidade}

Código, cadernos executados e registro de uso de IA:
\url{https://github.com/bsousa7/APPD_26}.
% Atualizar se o repositorio for renomeado para TrabalhoFinal_APPD_<sobrenomes>.
  % opcional: dicionario de variaveis, link do repositorio
\end{document}
